% Original proof note for integration into the research article.
% This is a fragment, without a document preamble.
% The optimization argument is self-contained. The general network-integrality
% principle is classical; no priority claim is made for that principle.

\section{Certified minimization of cocycle normal seeds}
\label{sec:cocycle}

\subsection{The precise family being optimized}

Let $\mathcal T$ be a finite triangulation of a compact orientable
$3$-manifold $M$. Generalized triangulations with distinct faces of the
same tetrahedron identified are allowed, provided they genuinely triangulate
the stated manifold. Ideal vertices must first be truncated. Choose an
orientation for every global edge, and let $c$ be an integral simplicial
$1$-cocycle. In tetrahedron $t$, write $v_{t,i}$ for the global vertex
represented by local corner $i\in\{0,1,2,3\}$. There are integral local heights
$h_{t,i}$, unique up to an additive constant for each tetrahedron, satisfying
\[
 h_{t,j}-h_{t,i}=c([i,j]_t).
\]
The cocycle equation on triangular faces guarantees their existence.
On a paired face the two lists of heights differ by a constant integer.

For an integral global vertex potential $f$, put
\[
 c_f=c+\delta f,\qquad a_{t,i}=h_{t,i}+f(v_{t,i}).
\]
The half-integral levels of the affine local height functions fit together
to give a cooriented normal surface $S(c_f)$. They fit because paired-face
height functions differ by an integer and we use all levels in
$\mathbb Z+\tfrac12$. Its signed intersection with each oriented edge is
$c_f(e)$; in particular it represents the relative homology class
Poincar\'e dual to $[c]$.

Order the four local corners so that
$a_0\leq a_1\leq a_2\leq a_3$. The three successive gaps count respectively
triangles at the lowest corner, quadrilaterals separating the lowest two
corners from the highest two, and triangles at the highest corner:
\[
 a_1-a_0,\quad a_2-a_1,\quad a_3-a_2.
\]
Zero gaps cause no ambiguity when heights coincide. Thus there is at most
one nonzero quadrilateral type, and the number of discs in tetrahedron $t$ is
exactly its height span. Our objective is consequently
\[
 D(f)=\sum_{t=1}^T\left(\max_i a_{t,i}-\min_i a_{t,i}\right).
 \tag{CS1}
\]
This is the number of normal discs, not the Euler characteristic, genus,
Thurston norm, or pattern-complexity. All coordinates are written in binary;
the construction takes three gap operations per tetrahedron.

\subsection{A sparse network and a short dual certificate}

\begin{theorem}[Integral span minimization]
Given $T$ lists of four integer heights and four global vertex labels per
list, the minimum of \textup{(CS1)} over integral potentials can be computed
in time polynomial in their binary input size. The same minimum is obtained
over real potentials. There is an optimality certificate consisting of an
integral potential and $T$ matched corner pairs. It can be checked with
$O(T)$ integer arithmetic operations after indexing the vertex labels.
\end{theorem}

\begin{proof}
Introduce free variables $L_t,U_t$ and $f_v$ and minimize
\[
 \sum_t(U_t-L_t)
\]
subject to
\[
 U_t-f_{v_{t,i}}\geq h_{t,i},\qquad
 f_{v_{t,i}}-L_t\geq-h_{t,i}
 \quad(t=1,\ldots,T;\ i=0,1,2,3).
 \tag{CS2}
\]
For fixed $f$, minimizing over $L,U$ recovers \textup{(CS1)} exactly.
Each constraint is a difference between two variables, so its coefficient
matrix is a directed incidence matrix, up to transposition. In particular,
this is an integral network linear program. The following explicit dual
description also yields a constructive proof of integrality.

Make nodes $\ell_t$, $u_t$, and one node for each global vertex $v$. Add
four arcs $\ell_t\to v_{t,i}$ of reward $-h_{t,i}$ and four arcs
$v_{t,i}\to u_t$ of reward $h_{t,i}$. Each $\ell_t$ has supply $1$, each
$u_t$ has demand $1$, and the vertex nodes have balance zero. Maximize
total reward. The network is feasible: send a unit from each $\ell_t$
through any one of its corners to $u_t$. Its underlying directed graph
has three layers, so a feasible flow has no directed circulation, and
every arc carries between zero and one unit. Network integrality gives an
integral optimum.

An integral flow selects one corner $(t,i)$ leaving each lower node and one
corner $(s,j)$ entering each upper node. At a global vertex, match the
incoming selected corners arbitrarily to the outgoing selected corners.
This produces $T$ rows $(t,i,s,j)$ such that sources $t$ and targets $s$
each occur once, and $v_{t,i}=v_{s,j}$ in every row. The dual value is
\[
 Q=\sum_{(t,i,s,j)}(h_{s,j}-h_{t,i}).
 \tag{CS3}
\]
For any real potential and any such row, the common vertex potential cancels:
\[
 \max_k a_{s,k}-\min_k a_{t,k}
 \geq a_{s,j}-a_{t,i}=h_{s,j}-h_{t,i}.
\]
Summing over the permutations of lower and upper nodes gives $D(f)\geq Q$.
Therefore a potential and a matching with $D(f)=Q$ certify optimality even
against all real potentials. This verifier needs no flow solver.

For completeness, an optimal flow gives an integral potential achieving
equality. In its residual graph retain every forward arc with its original
reward, and retain the reverse arc with negated reward whenever the original
flow is positive. Optimality says that there is no positive-reward directed
cycle. Add a universal source with zero-reward arcs and take longest-path
distances $x$. These are finite integers. For each residual arc $p\to q$
with reward $r$, they satisfy $x_q-x_p\geq r$. On each positive-flow arc
the reverse inequality is also present, giving equality. Set
$f_v=x_v$, $L_t=x_{\ell_t}$, and $U_t=x_{u_t}$. These satisfy
\textup{(CS2)}, and multiplying the equalities on used arcs by their flows
and summing gives
\[
 \sum_t(U_t-L_t)=Q.
\]
The actual span objective of this $f$ is at most the left side, while weak
duality gives the opposite inequality $D(f)\geq Q$. Hence equality holds.
\end{proof}

\paragraph{Implementation and binary complexity.}
The implementation negates rewards to obtain a min-cost problem, adds a
source and sink with unit source/sink capacities, and performs $T$
shortest-path augmentations. Each internal arc has capacity $T+1$; this is
strictly larger than the total flow, so it never removes an original
forward arc from the final residual graph. Parallel arcs from repeated
global vertices are retained; they encode different local heights.

There are $2T+V+2\leq6T+2$ nodes and $10T$ forward arcs. Bellman--Ford
uses $O(NE)$ relaxations per augmentation, giving $O(T^3)$ integer
relaxations overall, followed by $O(T^2)$ work to extract the potential.
If $H=\max(1,|h_{t,i}|)$, the rewards have magnitude at most $H$, and
simple-path distances have magnitude $O(TH)$. The resulting potential can
be translated by a constant so that its smallest entry is zero; all entries
still have magnitude $O(TH)$. Thus all arithmetic uses
$O(\log H+\log(T+1))$ bits. With ordinary binary addition/comparison, this
gives the conservative bit bound
\[
 O\bigl(T^3(\log H+\log(T+1))\bigr),
\]
up to vertex-indexing and input/output costs. The $T$ matching rows use
$O(T\log(T+1))$ bits; the potential and normal coordinates use
$O(T(\log H+\log(T+1)))$ bits. No capacity or iteration count depends on
the numerical value of a height or a normal coordinate.

\paragraph{A weighted variant.}
Positive integral weights $w_t$ on tetrahedral discs replace supply and
demand $1$ by $w_t$. The same network formulation and integrality proof
apply. A general polynomial-time min-cost flow algorithm handles binary
weights. The present code intentionally implements only the unweighted
case: replacing its $T$ unit augmentations by $\sum_t w_t$ unit
augmentations would be pseudopolynomial in binary weights.

\subsection{Edge coherence and deletion of zero classes}

Call a cooriented normal surface \emph{edge coherent} when all of its
intersections with any one oriented global edge have the same sign. An
edge with no intersections is allowed. Every surface $S(c_f)$ above is
edge coherent.

\begin{lemma}[Reconstruction from a coherent edge cocycle]
Let $S$ be an edge-coherent cooriented normal surface. Its signed edge
intersection numbers form an integral cocycle $c_S$, and its normal
coordinates equal those of $S(c_S)$.
\end{lemma}

\begin{proof}
The algebraic intersection count around a triangular face is zero: each
oriented normal arc contributes cancelling endpoints. Hence the signed
edge counts form a cocycle. Edge coherence implies that the unsigned
count on edge $e$ is $|c_S(e)|$.

Here is an explicit uniqueness argument for admissible normal coordinates
with specified unsigned edge counts. In a tetrahedron, write $x_i$ for its
four triangle counts and $q_{01},q_{02},q_{03}$ for its three quadrilateral
counts, indexed by the pair separated from its complement. The six edge
counts give a linear map from these seven coordinates. Its kernel is
one-dimensional, generated by
\[
 (1,1,1,1,-1,-1,-1).
\]
One can verify this directly: all four triangles together meet each edge
twice, as do all three quadrilaterals together; solving the six homogeneous
edge equations gives only this relation. Suppose two admissible coordinate
vectors had the same edge counts. Their difference is a multiple $\lambda$
of this vector. If $\lambda>0$, the second vector must have all three
quadrilateral entries at least $\lambda$, violating quadrilateral
compatibility; if $\lambda<0$, the same contradiction applies to the first
vector. Thus $\lambda=0$. Both $S$ and the local-height surface $S(c_S)$
are admissible and have the same unsigned edge counts, proving equality.
\end{proof}

\begin{theorem}[No zero-class component union at an optimum]
Let $f$ minimize \textup{(CS1)} for a cocycle $c$ in a compact orientable
$3$-manifold. No nonempty union of components of $S(c_f)$ represents zero
in $H_2(M,\partial M;\mathbb Z)$. In particular every component is
nonseparating within its ambient connected component of $M$.
\end{theorem}

\begin{proof}
Suppose a nonempty component union $Q$ represents zero. The complement
surface $S'=S(c_f)\setminus Q$ is still cooriented, normal, and edge
coherent, because deletion does not create a sign disagreement. The
signed edge cocycle of $Q$ represents its Poincar\'e-dual class, so it is
an integral coboundary $\delta g$. The remaining surface has cocycle
$c_f-\delta g=c+\delta(f-g)$. By the reconstruction lemma, its normal
coordinates are precisely those of $S(c+\delta(f-g))$.

Every nonempty normal surface contains at least one disc. Deletion of $Q$
therefore strictly reduces the objective \textup{(CS1)}, contradicting
minimality. Finally, a connected cooriented separating properly embedded
surface is the relative boundary of a component of its complement, so
it represents zero and is excluded by the result just proved.
\end{proof}

\begin{corollary}[Connected rank-one primitive seed]
Suppose $M$ is connected, $H^1(M;\mathbb Z)\cong\mathbb Z$, and $[c]$ is a
primitive generator. Every minimum-disc cocycle seed $S(c_f)$ is connected
and nonseparating.
\end{corollary}

\begin{proof}
Every component is nonseparating by the theorem. A connected cooriented
nonseparating surface represents a primitive relative homology class:
join a point on one side to a point on the other through its connected
complement and close across the surface. The resulting closed curve
has algebraic intersection number one with that surface. In rank one,
each component therefore represents either the positive or the negative
generator. Opposite signs would give a nonempty two-component union
with zero relative class, which the theorem excludes. All signs agree,
and their sum is the given primitive generator, so exactly one component
is present.
\end{proof}

\paragraph{Checking the hypotheses.}
The arbitrary-height API certifies only the optimization problem.
The face-pairing adapter checks reciprocal gluings, integer height
transitions, and normal matching; it does not certify manifold links.
The archive02 protocol adapter uses its separately validated finite
triangulation object. The connectedness theorem requires both rank one and
integral primitivity. In a spanning-tree gauge, integral cocycles form
the integral kernel of a matrix and hence a saturated sublattice. When
that kernel has rational dimension one, a nonzero kernel vector whose
coordinate gcd is one generates its integral lattice. This provides a
simple primitivity check for that specific rank-one construction.
An arbitrary rational cohomology basis in higher rank is not thereby a
saturated integral basis.

The initial spanning-tree gauge already has a related property: a
zero-class component union would give an exact edge cocycle vanishing on
every tree edge. Its vertex potential is then constant, so its cocycle
is zero; edge coherence forces the union to be empty. Thus the connected
rank-one seed is not a new existence theorem. The optimization result
shows that the guarantee survives leaving this special gauge, although
arbitrary changes of vertex potential can introduce separating components.

\subsection{What this supplies to a hierarchy implementation}

Theorem 6.2 of Lackenby's 2026 preprint \cite{lackenby2026} constructs nonseparating normal
surfaces from integral cocycles. The present routine supplies an exact
minimum-disc representative within the vertex-potential family of a
selected class, with an independent lower-bound certificate. The general
principle that suitable topological optimization problems become integral
linear programs is classical; for example, Dey, Hirani, and Krishnamoorthy
\cite{dey2011} treat optimal homologous cycles. The specific local-span objective, binary implementation,
and no-zero-component-union consequence above are the contribution being
documented here, without a claim of literature-wide priority.

There is a small normalization clarification in the proof of Lackenby's
Theorem 6.2, on page 27 of the PDF. After reducing a rational kernel to
one dimension, the condition that the sum of its coordinates is one need
not intersect that line: the line spanned by $(1,-1)$ is a simple example.
Set one known nonzero coordinate to one instead, then clear denominators.
This is an implementation-safe repair of that linear-algebra step; it
does not challenge the stated theorem. The normalizing row then has norm
one, so the existing Hadamard determinant estimate still applies.

\paragraph{Limits and further questions.}
The optimization does not minimize pattern-complexity or prove
incompressibility. A lower disc count is not a demonstrated reduction of
the hierarchy's number of resets. It supplies neither compressed cutting
with all attaching data nor a bound on hierarchy depth. In a one-vertex
triangulation, every local height receives the same global shift and the
objective is constant; the code returns a certificate of this limitation.
Truncation or subdivision can create useful vertex freedom, but increasing
the triangulation also changes the objective and costs work. No claim of
subdivision-invariant improvement or of completeness over all normal
surfaces is justified.

The genuine finite-exterior experiments therefore use the ten-vertex,
56-tetrahedron triangulations produced by truncating the two-tetrahedron
ideal trefoil and figure-eight exteriors. Simplifying those triangulations
first yields one-vertex triangulations and removes this particular freedom.
For the spanning-tree cocycles tested, the disc counts decrease from
109 to 54 and from 60 to 38. Independent Regina checks give connected
orientable surfaces of Euler characteristic $-1$ with one boundary
component. These are measurements of a normal-seed primitive on two
examples, not end-to-end recognition speedups or evidence of a general
asymptotic bound.

Useful next questions are whether the potential can be optimized jointly
with a small family of cohomology classes, whether subdivision choices
can be guided by a certified benefit/cost estimate, and whether disc
optimization can be combined with a separate certified pattern-complexity
descent. A weighted objective can favor expensive local handle types, but
its network algorithm must remain polynomial in binary capacities.

\subsection{A compressed Euler-characteristic check}

The normal coordinates also permit a useful positive-certificate test without
expanding the discs. Let $F=D(f)$, let $A$ be the sum of the numbers of sides
of all normal discs, and let $A_{\partial}$ be the total number of normal
arcs lying in unpaired boundary faces. If $q_{t,k}$ and $x_{t,i}$ are the
quadrilateral and triangle coordinates, then
\[
 A=3\sum_{t,i}x_{t,i}+4\sum_{t,k}q_{t,k}.
\]
Each internal normal arc is counted twice before gluing, and each boundary
arc once. Each point on a global triangulation edge is a vertex of the
surface cell structure. Thus
\begin{equation}
 \chi(S(c_f))=
 \sum_{e\text{ global}}|c_f(e)|
 -\frac{A+A_{\partial}}{2}+F.
 \label{eq:seed-euler}
\end{equation}
All terms use binary coordinates and a linear number of incidence records.
The formula assumes the verified manifold triangulation and normal matching
already required above; it is not a way to infer those hypotheses from an
arbitrary list of heights.

\begin{proposition}[A sufficient positive test]
Suppose $M$ has independently been identified as the exterior of the input
knot, and the primitive rank-one hypotheses for the optimized seed hold.
If the value in \eqref{eq:seed-euler} is one, the knot is trivial.
\end{proposition}
\begin{proof}
The seed is a connected orientable properly embedded surface. The
classification formula $\chi=2-2g-b$ shows that $\chi=1$ forces $g=0$ and
$b=1$, so it is a disk. Its nonzero primitive relative homology class implies
that its boundary is an essential longitude on the boundary torus of the
knot exterior. Joining that longitude to the core knot by the corresponding
annulus in its tubular neighborhood produces an embedded spanning disk.
\end{proof}

This is a proposed certificate interface, not a newly integrated recognition
stage. A value at most zero is inconclusive. Even for the unknot, minimizing
disc count within one cocycle's vertex-potential family has not been shown
to find a disk. In particular, the optimizer's globality is globality within
the stated family, not among all normal representatives of the homology class.
The package's product-solid-torus examples illustrate a successful case, while
the two knotted exteriors give $\chi=-1$.

\subsection{Transporting an optimum through known gauge changes}

The dual certificate has another direct performance consequence on a fixed
triangulation. Suppose new local heights are known to have the form
\[
 h'_{t,i}=h_{t,i}+g(v_{t,i})+a_t,
\]
where $g$ is a global integral potential and $a_t$ is an arbitrary integral
offset for tetrahedron $t$. These changes describe the same cocycle class.

\begin{proposition}[Certificate transport]
If $f$ and a matching certify the optimum for $h$, then $f'=f-g$ and the
same matching certify the same optimum for $h'$. The normal coordinates are
unchanged. Transport and verification take a linear number of binary
arithmetic operations after indexing.
\end{proposition}
\begin{proof}
The adjusted heights become
$h'_{t,i}+f'(v_{t,i})=h_{t,i}+f(v_{t,i})+a_t$. Every local span and gap is
unchanged. On a matched row $(t,i,s,j)$, the common global vertex cancels
the added $g$ values, leaving an added reward $a_s-a_t$. Sources and targets
are each permutations of the tetrahedra, so these offsets sum to zero.
The old primal/dual equality therefore remains valid. The reported initial
disc count must of course be recomputed for the new input before serializing
a certificate.
\end{proof}

This transport formula is a proved opportunity for a future cache interface;
the patch does not expose a new transport API. It also sharpens the
interpretation of the large-coboundary experiment: if the generating
potential and an old optimum are supplied, one can transport that optimum
without solving a new flow. The experiment deliberately asks the general
optimizer to recover an optimum from the changed heights alone. Cutting or
retriangulating changes the incidence data and is not covered by this formula.
